\chapter{Lagrangeova dualnost}
\label{ch:dualnost}

\begin{izvor}
\textbf{Izvori:} \texttt{04\_Lagrangeova\_dualnost.pdf} (str.\ 1--34) i
\texttt{Vjezbe\_3.pdf} (str.\ 1--42) --- oboje obrađeno u cijelosti.

Vježbe 3 imaju dva primjera. \textbf{Primjer 1} (tržište i formiranje cijena)
je primjena dualne dekompozicije, pa stoji uz nju
(odjeljak~\ref{sec:dekompozicija}). \textbf{Primjer 2} dodiruje više tema, pa je
razlomljen po podzadacima: a) i b) uz ponovljene KKT uvjete
(odjeljak~\ref{sec:kkt-opet}), c) i d) uz dualni problem
(odjeljak~\ref{sec:primjer2-dual}), e) i f) uz analizu osjetljivosti
(odjeljak~\ref{sec:osjetljivost}). Oznaka \str{V3:20} znači „Vježbe 3,
stranica 20”.
\end{izvor}

\textbf{Što ovo poglavlje donosi.} U poglavlju~\ref{ch:uvjeti} Lagrangeova
funkcija bila je samo \emph{tehnički trik} za izvod nužnih uvjeta. Ovdje ista
funkcija dobiva dublje značenje: iz nje se gradi \textbf{drugi optimizacijski
problem} (dualni), koji daje donju granicu na rješenje polaznog i --- pod
određenim uvjetima --- \emph{istu} optimalnu vrijednost. To otvara tri stvari:

\begin{enumerate}
  \item \textbf{provjeru optimalnosti} („koliko sam daleko od optimuma?”),
  \item \textbf{distribuirano rješavanje} (veliki problem rastavljen na male,
        neovisne),
  \item \textbf{analizu osjetljivosti} (koliko se optimum mijenja ako popustimo
        ograničenje).
\end{enumerate}

% =====================================================================
\section{Donje granice i dualna funkcija}
\label{sec:donje-granice}
\str{4--9}

\textbf{Optimizacijski problem} \str{4}:
\begin{equation}
\begin{aligned}
\minim & f(x) \\
\uzogr & h_i(x) = 0, && i = 1,\dots,m \\
       & g_j(x) \le 0, && j = 1,\dots,p
\end{aligned}
\end{equation}

Neka je $p\opt$ \textbf{optimalna vrijednost} $f(x)$.

\subsection{Ključno opažanje}
\str{5}

\begin{teorem}[Donje granice]
Neka je $x$ u dozvoljenom skupu ($g(x) \le 0$, $h(x) = 0$). Za proizvoljan
$\lambda \in \R^m$ i $\mu \in \R^p$ gdje je $\mu \ge 0$ imamo
\begin{equation}
\label{eq:donja-granica}
L(x,\lambda,\mu) := f(x) + \sum_{i=1}^{m}\lambda_i h_i(x)
                        + \sum_{j=1}^{p}\mu_j g_j(x) \;\le\; f(x).
\end{equation}
\end{teorem}

\textbf{Zašto to vrijedi --- provjerite član po član:}

\begin{itemize}
  \item $\sum_i \lambda_i h_i(x)$: budući da je $x$ dopustiv, vrijedi
        $h_i(x) = 0$ za svaki $i$, pa je cijela ta suma \textbf{jednaka nuli},
        bez obzira na $\lambda$. Zato $\lambda$ smije biti bilo kojeg predznaka.
  \item $\sum_j \mu_j g_j(x)$: dopustivost daje $g_j(x) \le 0$, a
        pretpostavka daje $\mu_j \ge 0$. Umnožak nenegativnog i nepozitivnog
        broja je \textbf{nepozitivan}, pa je i cijela suma $\le 0$.
  \item Dakle $L = f + 0 + (\text{nešto} \le 0) \le f$. \checkmark
\end{itemize}

\begin{napomena}[ovdje se vidi odakle uvjet $\mu \ge 0$]
Kad bi neki $\mu_j$ bio negativan, član $\mu_jg_j(x)$ mogao bi biti
\emph{pozitivan} i cijela bi nejednakost pala. Zahtjev $\mu \ge 0$ je upravo
ono što čini $L$ donjom ogradom za $f$ na dozvoljenom skupu.
\end{napomena}

\subsection{Infimizacija obiju strana}
\str{6}

Nakon infimizacije lijeve i desne strane \eqref{eq:donja-granica}:
\begin{equation}
\label{eq:lanac}
\ell(\lambda,\mu) := \inf_x L(x,\lambda,\mu)
\;\le\; \inf_{\{x \mid g(x)\le0,\,h(x)=0\}} L(x,\lambda,\mu)
\;\le\; \inf_{\{x \mid g(x)\le0,\,h(x)=0\}} f(x) = p\opt
\end{equation}

\textbf{Objašnjenje svakog koraka u lancu} --- to je najvažniji izvod poglavlja:

\begin{enumerate}
  \item \textbf{Prva nejednakost.} Lijevo se infimum uzima po \emph{svim}
        $x \in \R^n$, desno samo po \emph{dopustivim} $x$. Traženje minimuma po
        \textbf{većem} skupu daje \textbf{manju ili jednaku} vrijednost --- jer
        veći skup uključuje sve kandidate manjeg, i možda još bolje.
  \item \textbf{Druga nejednakost.} Ovo je izravno \eqref{eq:donja-granica}
        („$L \le f$ na dopustivim točkama”), primijenjeno pod istim infimumom.
  \item \textbf{Zadnja jednakost.} Infimum od $f$ po dopustivim točkama je po
        definiciji optimalna vrijednost $p\opt$.
\end{enumerate}

\textbf{Rezultat:} $\ell(\lambda,\mu) \le p\opt$ za \emph{svaki} par
$(\lambda,\mu)$ s $\mu \ge 0$. Dobili smo \textbf{stroj za proizvodnju donjih
granica}: uvrstite bilo koje dopušteno $(\lambda,\mu)$ i dobit ćete broj za
koji sigurno znate da nije veći od pravog optimuma.

\subsection{Najbolja donja granica}
\str{9}

Budući da su $\lambda$ i $\mu \ge 0$ proizvoljni, prirodno je tražiti
\textbf{najbolju} (najveću) takvu granicu:
\begin{equation}
\label{eq:slaba-dualnost}
d\opt = \max_{\{\lambda,\mu \mid \mu\ge0\}} \ell(\lambda,\mu)
\;\le\; \inf_{\{x \mid g(x)\le0,\,h(x)=0\}} f(x) = p\opt
\end{equation}

% =====================================================================
\section{Terminologija}
\str{10}

\begin{definicija}[Pojmovi dualnosti]
\begin{itemize}
  \item \textbf{Lagrangian} (Lagrangeova funkcija):
        \[
        L(x,\lambda,\mu) := f(x) + \sum_{i=1}^{m}\lambda_ih_i(x)
                                 + \sum_{j=1}^{p}\mu_jg_j(x)
        \]
  \item \textbf{Lagrangeova dualna funkcija cilja:}
        $\ell(\lambda,\mu) := \inf_x L(x,\lambda,\mu)$
  \item \textbf{Lagrangeov dualni problem:}
        $d\opt = \sup_{\{\lambda,\mu \mid \mu\ge0\}} \ell(\lambda,\mu)$
  \item \textbf{Primarni problem:}
        $p\opt = \inf_{\{x \mid g(x)\le0,\,h(x)=0\}} f(x)$
  \item $\lambda_i$ je Lagrangeov množitelj (\textbf{dualna varijabla})
        pridružen ograničenju $h_i(x) = 0$
  \item $\mu_j$ je Lagrangeov množitelj (dualna varijabla) pridružen
        ograničenju $g_j(x) \le 0$
\end{itemize}
\end{definicija}

\begin{napomena}[dvije razine minimizacije --- ne pobrkati ih]
\begin{itemize}
  \item \textbf{Unutarnja} operacija je $\inf_x$: minimizira se po
        \emph{primarnim} varijablama $x$, uz $(\lambda,\mu)$ \emph{fiksne}.
        Rezultat je broj koji ovisi o $(\lambda,\mu)$ --- to je $\ell$.
  \item \textbf{Vanjska} operacija je $\max$: maksimizira se po
        \emph{dualnim} varijablama.
\end{itemize}
Dakle dualni problem je \emph{max-min}, dok je primarni obični \emph{min}.
Uočite i da unutarnji $\inf_x$ ide po \textbf{cijelom} $\R^n$ --- ograničenja
su „progutana” u $L$.
\end{napomena}

% =====================================================================
\section{Slaba dualnost}
\str{11}

\begin{teorem}[Slaba dualnost --- donje granice]
\[
d\opt = \max_{\{\lambda,\mu \mid \mu\ge0\}} \ell(\lambda,\mu)
\;\le\; \inf_{\{x \mid g(x)\le0,\,h(x)=0\}} f(x) = p\opt
\]
\textbf{Dualna optimalna vrijednost} $(d\opt)$ $\le$ \textbf{Primarna optimalna
vrijednost} $(p\opt)$.

\textbf{Slaba dualnost uvijek vrijedi!} (Bez obzira je li primarni problem
konveksan ili ne.)
\end{teorem}

\textbf{Druga ključna tvrdnja sa slajda:} \textbf{Dualni problem je konkavan
maksimizacijski problem.}

\begin{napomena}[zašto je dualni problem uvijek konkavan]
Za \emph{fiksan} $x$, funkcija $(\lambda,\mu) \mapsto L(x,\lambda,\mu)$ je
\textbf{afina} u $(\lambda,\mu)$ --- multiplikatori ulaze samo linearno. Dualna
funkcija $\ell$ je infimum \emph{cijele obitelji} takvih afinih funkcija (po
svim $x$), a infimum afinih funkcija je uvijek \textbf{konkavna} funkcija.
Ograničenje $\mu \ge 0$ je afino.

\textbf{Posljedica:} dualni problem je \emph{uvijek} konveksan (maksimizacija
konkavne funkcije), čak i kad je primarni problem grozno nekonveksan. Zato je
dualnost tako moćna --- vidi Primjer 2 niže.
\end{napomena}

\textbf{Dualni jaz.} Razlika $p\opt - d\opt \ge 0$ zove se \textbf{dualni jaz}
(engl.\ \emph{duality gap}). Slaba dualnost kaže da je uvijek nenegativan.
Praktična korist: ako riješimo dualni problem i dobijemo $d\opt$, a imamo bilo
koju dopustivu točku $\bar x$, onda znamo
\[
d\opt \le p\opt \le f(\bar x)
\]
--- dakle znamo \emph{koliko najviše} možemo još dobiti. To je „certifikat
kvalitete” rješenja.

% =====================================================================
\section{Primjer 1: linearno programiranje u standardnoj formi}
\label{sec:dual-lp}
\str{12--13}

\begin{equation}
\begin{aligned}
\minim & c\T x \\
\uzogr & Ax - b = 0, \\
       & x \ge 0 \qquad (-x \le 0)
\end{aligned}
\end{equation}

\textbf{Lagrangian:}
\[
L(x,\lambda,\mu) = c\T x + \lambda\T(Ax-b) - \mu\T x
= -b\T\lambda + (c + A\T\lambda - \mu)\T x
\]

\emph{Raspis preslagivanja (na slajdu nije prikazan):}
\begin{align}
L &= c\T x + \lambda\T Ax - \lambda\T b - \mu\T x \\
  &= -\lambda\T b + \big(c\T + \lambda\T A - \mu\T\big)x \\
  &= -b\T\lambda + \big(c + A\T\lambda - \mu\big)\T x
\end{align}
\emph{Međukoraci:} $\lambda\T b = b\T\lambda$ (skalar je jednak svom
transponatu); $\lambda\T A = (A\T\lambda)\T$; skupljeni koeficijenti uz $x$ daju
vektor $c + A\T\lambda - \mu$.

\textbf{Uočite:} $L$ je \textbf{afina funkcija u $x$}. Stoga imamo
\begin{equation}
\ell(\lambda,\mu) = \inf_x L(x,\lambda,\mu) =
\begin{cases}
-b\T\lambda & \text{ako je } c + A\T\lambda - \mu = 0 \\
-\infty & \text{u ostalim slučajevima}
\end{cases}
\end{equation}

\begin{napomena}[zašto afina funkcija ima infimum $-\infty$]
Afina funkcija $\alpha + \beta\T x$ nema donju granicu osim ako je
$\beta = 0$: uzmemo li $x = -t\beta$ za veliki $t$, vrijednost ide u
$-\infty$. Ako je pak $\beta = 0$, funkcija je konstanta $\alpha$ i infimum je
upravo $\alpha$.

Ovdje je $\beta = c + A\T\lambda - \mu$ i $\alpha = -b\T\lambda$ --- odatle
gornji rezultat. Vrijednost $-\infty$ nije problem: takav $(\lambda,\mu)$ daje
beskorisnu donju granicu, pa ga maksimizacija sama odbacuje.
\end{napomena}

\textbf{Donja granica:} $p\opt \ge -b\T\lambda$ za bilo koji $\lambda$ takav da
je $A\T\lambda + c \ge 0$.

\emph{Odakle taj uvjet:} iz $c + A\T\lambda - \mu = 0$ slijedi
$\mu = c + A\T\lambda$; a budući da mora biti $\mu \ge 0$, dobivamo
$A\T\lambda + c \ge 0$. Time je $\mu$ eliminiran iz zapisa.

\textbf{Najviša (najbolja) donja granica:}
\[
p\opt \ge \max_\lambda\{\, -b\T\lambda \mid A\T\lambda + c \ge 0 \,\}
\]

\textbf{Poanta:} dual jednog LP-a je opet LP. To je klasična „dualnost LP-a”
koju ćemo koristiti u poglavljima~\ref{ch:lpqp}~i~\ref{ch:robusno}.

% =====================================================================
\section{Primjer 2: optimalno particioniranje u dva skupa}
\label{sec:particioniranje}
\str{14--15}

\textbf{Postavka} \str{14}\textbf{.} Neka je $W \in \R^{n\times n}$ simetrična
matrica:
\begin{equation}
\begin{aligned}
\minim & x\T W x \\
\uzogr & x_i^2 = 1, \quad i = 1,\dots,n
\end{aligned}
\end{equation}

\textbf{Interpretacija (sa slajda):} elemente indeksirane skupom
$\{1,\dots,n\}$ potrebno je optimalno podijeliti u dva skupa koji su
enkodirani s $x_i = 1$ i $x_i = -1$.

\begin{itemize}
  \item $W_{ij}$ (element matrice $W$ u $i$-tom retku i $j$-tom stupcu) je
        \textbf{cijena stavljanja $i$ i $j$ u isti skup};
  \item $-W_{ij}$ je cijena stavljanja tih elemenata u odvojene skupove.
\end{itemize}

\emph{Zašto to radi:} $x\T Wx = \sum_{i,j} W_{ij}x_ix_j$. Ako su $i$ i $j$ u
istom skupu, $x_ix_j = +1$ i pribraja se $W_{ij}$; ako su u različitim,
$x_ix_j = -1$ i pribraja se $-W_{ij}$. Uvjet $x_i^2 = 1$ znači točno
$x_i \in \{-1, +1\}$.

\textbf{Tvrdnja sa slajda:} \emph{Ovo je nekonveksan problem. Dozvoljeni skup
sadrži $2^n$ diskretnih točaka.}

\subsection{Dualna funkcija}
\str{15}

Uz definiciju $\mathbf{1} = \begin{pmatrix}1&1&\dots&1\end{pmatrix}\T$:
\begin{align}
\ell(\lambda)
&= \inf_x\left(x\T Wx + \sum_{i=1}^{n}\lambda_i(x_i^2 - 1)\right) \\
&= \inf_x\; x\T\big(W + \diag(\lambda)\big)x - \mathbf{1}\T\lambda
\end{align}

\emph{Raspis:} $\sum_i \lambda_i x_i^2 = x\T\diag(\lambda)x$, gdje je
$\diag(\lambda)$ dijagonalna matrica s $\lambda_i$ na dijagonali. Preostali dio
$-\sum_i\lambda_i = -\mathbf{1}\T\lambda$ ne ovisi o $x$, pa izlazi ispred
infimuma. \checkmark

\begin{equation}
\ell(\lambda) =
\begin{cases}
-\mathbf{1}\T\lambda & \text{ako je } W + \diag(\lambda) \succeq 0 \\
-\infty & \text{u ostalim slučajevima}
\end{cases}
\end{equation}

\emph{Zašto:} kvadratna forma $x\T Mx$ ima infimum $0$ (postignut u $x = 0$)
ako i samo ako je $M \succeq 0$; inače postoji smjer $x$ u kojem je forma
negativna, pa skaliranjem $x \to tx$ ide u $-\infty$ (jer je forma kvadratna u
$t$). Ovo je izravna primjena definicije pozitivne semidefinitnosti
\str{2:46}.

\textbf{Donja granica:} $p\opt \ge -\mathbf{1}\T\lambda$ za bilo koji $\lambda$
takav da je $W + \diag(\lambda) \succeq 0$.

\textbf{Najviša donja granica:}
\[
p\opt \ge \max_\lambda\{\, -\mathbf{1}\T\lambda \mid
W + \diag(\lambda) \succeq 0 \,\}
\qquad \textbf{konveksan opt.\ problem!}
\]

\begin{napomena}[zašto je ovaj primjer spektakularan]
Polazni problem je \textbf{kombinatoran}: $2^n$ diskretnih točaka, dakle za
$n = 100$ oko $10^{30}$ mogućnosti --- nerješivo iscrpnim pretraživanjem.

Njegov \textbf{dual} je gladak konveksan problem s ograničenjem
$W + \diag(\lambda) \succeq 0$, tzv.\ \textbf{linearnom matričnom
nejednadžbom}. To je \emph{semidefinitni program} (SDP,
poglavlje~\ref{ch:qcpsdp}), koji se rješava učinkovito.

Nema jamstva da je $d\opt = p\opt$ (problem nije konveksan, pa jaka dualnost ne
mora vrijediti), ali dobivamo \textbf{netrivijalnu donju granicu} za problem
koji inače ne bismo mogli ni približno riješiti. Ovo je poznata SDP-relaksacija
problema \emph{max-cut}.
\end{napomena}

% =====================================================================
\section{Jaka dualnost i Slaterovi uvjeti}
\label{sec:jaka-dualnost}
\str{16--18}

\begin{teorem}[Teorem Lagrangeove dualnosti \str{16}]
\begin{itemize}
  \item Slaba dualnost uvijek vrijedi: $d\opt \le p\opt$.
  \item Kad je primarni problem \textbf{konveksan} i kad su zadovoljeni
        \textbf{Slaterovi uvjeti}, vrijedi \textbf{jaka dualnost}:
        $d\opt = p\opt$.
\end{itemize}

\textbf{Jaka dualnost u kompaktnoj formi:}
\[
\max_{\{\lambda,\mu \mid \mu\ge0\}}
\left(\inf_x\; f(x) + \lambda\T h(x) + \mu\T g(x)\right)
= \inf_{\{x \mid g(x)\le0,\,h(x)=0\}} f(x)
\]
\end{teorem}

\textbf{Riječima:} kod konveksnog problema smijemo \textbf{zamijeniti redoslijed}
operacija $\max$ i $\inf$ bez gubitka. To je vrlo netrivijalna tvrdnja ---
općenito je uvijek $\max\inf \le \inf\max$, a jednakost je iznimka.

\begin{definicija}[Slaterovi uvjeti (engl.\ \emph{Slater's constraint
qualification}) \str{17}]
Definirajmo skupove $I_n$, $I_a$: $i \in I_n$ ako je $g_i(\cdot)$
\textbf{nelinearna} funkcija; $i \in I_a$ ako je $g_i(\cdot)$ \textbf{afina}
funkcija.

\textbf{Slaterovi uvjeti:} skup
\[
\{\, x \mid h(x) = 0,\; g_i(x) < 0 \text{ za } i \in I_n,\;
g_i(x) \le 0 \text{ za } i \in I_a \,\}
\]
nije prazan skup.
\end{definicija}

\textbf{Kako to čitati.} Traži se postojanje barem jedne točke koja:

\begin{itemize}
  \item zadovoljava sve jednakosti \emph{točno};
  \item zadovoljava sva \textbf{nelinearna} ograničenja nejednakosti
        \textbf{strogo} (sa $<$, dakle \emph{s rezervom}, ne na samom rubu);
  \item zadovoljava \textbf{afina} ograničenja obično (dovoljno $\le$).
\end{itemize}

Takva se točka zove \textbf{strogo dopustiva}. Uvjet neformalno kaže da
dozvoljeni skup nije „spljošten” --- da ima nepraznu unutrašnjost u nelinearnim
smjerovima.

\begin{napomena}[praktično značenje]
Slaterovi uvjeti gotovo su uvijek ispunjeni u inženjerskim problemima; obično se
provjere jednom rečenicom („postoji točka koja strogo zadovoljava sva
nelinearna ograničenja”). Ako su \emph{sva} ograničenja nejednakosti afina
(npr.\ kod LP-a i QP-a s linearnim ograničenjima), uvjet se svodi na to da
dozvoljeni skup nije prazan --- dakle praktički besplatno.
\end{napomena}

\subsection{Primjer: LP u standardnoj formi, sada s jakom dualnošću}
\str{18}

Za problem iz odjeljka~\ref{sec:dual-lp}, dualna funkcija je
\[
\ell(\lambda,\mu) = \inf_x L(x,\lambda,\mu) =
\begin{cases}
-b\T\lambda & c + A\T\lambda - \mu = 0 \\
-\infty & \text{u ostalim slučajevima}
\end{cases}
\]

Donja granica: $p\opt \ge -b\T\lambda$ za bilo koji $\lambda$ takav da je
$A\T\lambda + c \ge 0$. \textbf{Zbog jake dualnosti:}
\[
p\opt = d\opt = \max_\lambda\{\, -b\T\lambda \mid A\T\lambda + c \ge 0 \,\}
\]

Nejednakost je postala \textbf{jednakost} --- LP je konveksan i (uz nepraznu
dopustivost) zadovoljava Slaterove uvjete, pa dualni problem daje \emph{točno}
istu vrijednost.

% =====================================================================
\section{KKT uvjeti (još jednom)}
\label{sec:kkt-opet}
\str{19--23}

Sada iz dualnosti \emph{ponovno} izvodimo KKT uvjete --- ali ovaj put dobivamo i
odgovor na pitanje kada su oni \textbf{dovoljni}.

\subsection{Jaka dualnost i uvjeti komplementarnosti}
\str{20--21}

\textbf{Pretpostavimo da vrijedi jaka dualnost.} Neka je $x\opt$ optimalan za
primarni problem, a $(\lambda\opt, \mu\opt)$ za dualni problem. Tada vrijedi:
\begin{align}
f(x\opt) = \ell(\lambda\opt,\mu\opt)
&= \inf_x\left(f(x) + \sum_{i=1}^{m}\lambda_i\opt h_i(x)
              + \sum_{j=1}^{p}\mu_j\opt g_j(x)\right)
  \label{eq:kkt-lanac1}\\
&\le f(x\opt) + \sum_{i=1}^{m}\lambda_i\opt h_i(x\opt)
              + \sum_{j=1}^{p}\mu_j\opt g_j(x\opt)
  \label{eq:kkt-lanac2}\\
&\le f(x\opt)
  \label{eq:kkt-lanac3}
\end{align}

\textbf{Objašnjenje svakog koraka:}

\begin{itemize}
  \item prva jednakost je \emph{jaka dualnost} ($p\opt = d\opt$) plus definicija
        dualne funkcije;
  \item \eqref{eq:kkt-lanac2}: infimum po \emph{svim} $x$ je $\le$ vrijednosti u
        \emph{jednoj konkretnoj} točki $x\opt$;
  \item \eqref{eq:kkt-lanac3}: to je nejednakost \eqref{eq:donja-granica} --- na
        dopustivoj točki dodani članovi su $\le 0$.
\end{itemize}

\textbf{Zaključak sa slajda:} zadnje dvije nejednakosti u biti vrijede
\textbf{sa znakom jednakosti} --- jer lanac počinje i završava istim brojem
$f(x\opt)$, pa ništa „u sredini” ne smije biti strogo manje:
\[
f(x\opt) = \ell(\lambda\opt,\mu\opt)
= \inf_x\Big(f(x) + \textstyle\sum_i\lambda_i\opt h_i(x)
                  + \sum_j\mu_j\opt g_j(x)\Big)
= f(x\opt) + \textstyle\sum_i\lambda_i\opt h_i(x\opt)
           + \sum_j\mu_j\opt g_j(x\opt)
\]

\textbf{Iz toga slijede dvije stvari} \str{21}:

\begin{enumerate}
  \item \textbf{$x\opt$ minimizira $L(x,\lambda\opt,\mu\opt)$.}

        \emph{Zašto:} jednakost u \eqref{eq:kkt-lanac2} znači da je infimum
        Lagrangiana postignut baš u $x\opt$.
  \item \textbf{$\mu_j\opt g_j(x\opt) = 0$} za sve $j = 1,\dots,p$ ---
        \textbf{uvjeti komplementarnosti} (već smo ih susretali!):
        \[
        g_j(x\opt) < 0 \Rightarrow \mu_j\opt = 0, \qquad
        \mu_j\opt > 0 \Rightarrow g_j(x\opt) = 0
        \]

        \emph{Zašto:} jednakost u \eqref{eq:kkt-lanac3} znači da je
        $\sum_j\mu_j\opt g_j(x\opt) = 0$ (članovi s $h_i$ su ionako nula). Ali
        to je zbroj \emph{nepozitivnih} članova; zbroj nepozitivnih brojeva je
        nula samo ako je \emph{svaki} od njih nula. \checkmark
\end{enumerate}

\subsection{Teorem: KKT iz dualnosti}
\str{22}

\begin{teorem}[KKT uvjeti]
Ako vrijedi jaka dualnost i $x\opt, \lambda\opt, \mu\opt$ su optimalni, oni
zadovoljavaju KKT uvjete. Prvi od njih je stacionarnost Lagrangeove funkcije,
$\nabla_x L(x\opt,\lambda\opt,\mu\opt) = 0$, raspisana po članovima:
\begin{align}
\nabla f(x\opt) + \sum_{i=1}^{m}\lambda_i\opt\nabla h_i(x\opt)
                + \sum_{j=1}^{p}\mu_j\opt\nabla g_j(x\opt) &= 0
  && \text{(grad.\ jed.)} \notag\\
h_i(x\opt) &= 0, && \text{(primarna ogr.)} \notag\\
g_j(x\opt) &\le 0, && \text{(primarna ogr.)} \notag\\
\mu_j\opt g_j(x\opt) &= 0, && \text{(komplementarnost)} \notag\\
\mu_j\opt &\ge 0, && \text{(dualna ogr.)} \notag
\end{align}
\end{teorem}

\emph{Odakle gradijentna jednadžba:} iz tvrdnje 1 gore, $x\opt$ minimizira
glatku funkciju $L(\cdot,\lambda\opt,\mu\opt)$ \emph{bez ograničenja}, pa je u
njoj gradijent nula --- uvjet iz odjeljka~\ref{sec:uvjeti-bez-ogr}.

\subsection{KKT uvjeti za konveksne probleme --- sada i dovoljni}
\str{23}

\begin{teorem}[Dovoljnost]
Ako $x\opt, \lambda\opt, \mu\opt$ zadovoljavaju KKT uvjete za \textbf{konveksan}
problem, onda su \textbf{optimalni}.
\end{teorem}

\textbf{Dokaz (sa slajda, s dopunjenim objašnjenjima).} Pretpostavimo da
$x\opt, \lambda\opt, \mu\opt$ zadovoljavaju KKT uvjete. Imamo da:

\begin{itemize}
  \item \textbf{iz komplementarnosti slijedi} $f(x\opt) =
        L(x\opt,\lambda\opt,\mu\opt)$.

        \emph{Zašto:} $L = f + \sum\lambda_ih_i + \sum\mu_jg_j$; prva suma je
        nula jer je $h_i(x\opt) = 0$, druga je nula po komplementarnosti.
  \item \textbf{iz gradijentne jednakosti i konveksnosti} ($f$, $h_i$ i $g_j$
        su sve konveksne funkcije) \textbf{slijedi}
        $\ell(\lambda\opt,\mu\opt) = L(x\opt,\lambda\opt,\mu\opt)$.

        \emph{Zašto:} uz $\mu\opt \ge 0$ i konveksnost, funkcija
        $x \mapsto L(x,\lambda\opt,\mu\opt)$ je konveksna (zbroj konveksnih uz
        nenegativne težine). Kod konveksne funkcije stacionarna točka je
        \emph{globalni} minimum (odjeljak~\ref{sec:svojstva-konv}), pa je
        infimum postignut u $x\opt$ --- a taj infimum je po definiciji $\ell$.
\end{itemize}

pa zaključujemo da je $f(x\opt) = \ell(\lambda\opt,\mu\opt)$, a budući da
$f(x\opt)$ ne može biti manji od $\ell(\lambda\opt,\mu\opt)$ (slaba dualnost),
nalazimo se nužno u optimumu. $\blacksquare$

\begin{teorem}[Karakterizacija optimuma \str{23}]
Ako imamo konveksan optimizacijski problem sa zadovoljenim Slaterovim uvjetima,
onda je $x\opt$ optimalan \textbf{ako i samo ako} postoje $\lambda\opt,
\mu\opt$ koji zadovoljavaju KKT uvjete.
\end{teorem}

\begin{napomena}[ovo je centralni praktični rezultat kolegija]
Za konveksan problem sa Slaterovim uvjetima, \textbf{rješavanje KKT sustava je
ekvivalentno rješavanju problema}. Nema više „kandidata” koje treba dodatno
provjeravati (usporedi upozorenja u poglavlju~\ref{ch:uvjeti}): svako rješenje
KKT sustava \emph{jest} globalni optimum.
\end{napomena}

\subsection{Vježbe 3, Primjer 2, podzadaci a) i b)}
\label{sec:primjer2-ab}

\begin{zadatak}[Vježbe 3, Primjer 2: dualnost, KKT, osjetljivost]
\label{zad:primjer2}
\str{V3:20}

\textbf{Postavka.}
\[
\begin{aligned}
\text{minimiziraj}_x \quad & f(x) = (x_1-3)^2 + (x_2-2)^2 \\
\text{uz ogr.} \quad
& -x_1 \le 0 && (g_1) \\
& x_1 - 2 \le 0 && (g_2) \\
& -x_2 \le 0 && (g_3) \\
& x_2 - 1 \le 0 && (g_4)
\end{aligned}
\]

\textbf{Podzadaci:}
\begin{enumerate}[label=\alph*)]
  \item Je li ovo konveksan opt.\ problem?
  \item Riješi opt.\ problem direktno rješavajući KKT uvjete.
  \item Riješi problem u MATLAB-u (\texttt{quadprog} i/ili YALMIP) --- odredi
        primarne i dualne varijable u optimumu.
  \item Postavi i riješi dualni problem (koristite MATLAB za rješavanje).
  \item Što možete reći o osjetljivosti optimalne vrijednosti ($p\opt =$ opt.\
        vrijednost primarnog problema) na perturbacije u ograničenjima?
  \item Ponovite (d) i (e), ali sad s promijenjenim ograničenjima $g_2$ i $g_4$
        ($g_1$ i $g_3$ ostaju ista):
        $g_2(x) = x_1 - 1{,}5 \le 0$, $g_4(x) = x_2 - 1{,}5 \le 0$.
\end{enumerate}

\medskip
\textbf{Geometrija problema.} Funkcija cilja je kvadrat udaljenosti od točke
$(3,2)$. Ograničenja definiraju \textbf{pravokutnik}
$0 \le x_1 \le 2$, $0 \le x_2 \le 1$. Točka $(3,2)$ leži \emph{izvan} tog
pravokutnika (gore desno), pa se traži točka pravokutnika najbliža njoj.

\medskip
\textbf{a) Je li ovo konveksan opt.\ problem?} \str{V3:21--22}

\begin{itemize}
  \item Funkcije ograničenja nejednakosti su konveksne (ovdje su one
        \textbf{afine}, pa stoga i konveksne).
  \item \textbf{Hessian:}
        \[
        H = \begin{bmatrix}
        \frac{\partial^2 f}{\partial x_1^2} &
        \frac{\partial^2 f}{\partial x_1\partial x_2} \\[4pt]
        \frac{\partial^2 f}{\partial x_2\partial x_1} &
        \frac{\partial^2 f}{\partial x_2^2}
        \end{bmatrix}
        = \begin{pmatrix} 2 & 0 \\ 0 & 2 \end{pmatrix} \succ 0
        \]
\end{itemize}

\emph{Raspis Hessiana:} iz $f = (x_1-3)^2 + (x_2-2)^2$ slijedi
$\partial f/\partial x_1 = 2(x_1-3)$, pa
$\partial^2 f/\partial x_1^2 = 2$; mješovita derivacija je $0$ jer prvi član ne
ovisi o $x_2$. Analogno za drugu varijablu. Matrica $2I$ ima obje svojstvene
vrijednosti jednake $2 > 0$, dakle je pozitivno definitna. \checkmark

\textbf{Zaključak sa slajda:} \emph{Konveksan optimizacijski problem sa
zadovoljenim Slaterovim uvjetima} $\to$ \textbf{KKT uvjeti su nužni i dovoljni
za optimalnost} $\to$ \textbf{vrijedi jaka dualnost}.

\emph{Provjera Slatera:} sva su ograničenja afina, pa je dovoljno da dozvoljeni
skup nije prazan --- npr.\ točka $(1;\ 0{,}5)$ zadovoljava sva četiri.
\checkmark

\medskip
\textbf{b) Rješavanje KKT uvjeta} \str{V3:23--27}

\textbf{Lagrangian:}
\[
L(x,\mu) = (x_1-3)^2 + (x_2-2)^2 + \mu_1(-x_1) + \mu_2(x_1-2)
         + \mu_3(-x_2) + \mu_4(x_2-1)
\]

\textbf{KKT uvjeti:}
\begin{align}
\begin{pmatrix} 2x_1 - 6 \\ 2x_2 - 4 \end{pmatrix}
+ \begin{pmatrix} -1 \\ 0 \end{pmatrix}\mu_1
+ \begin{pmatrix} 1 \\ 0 \end{pmatrix}\mu_2
+ \begin{pmatrix} 0 \\ -1 \end{pmatrix}\mu_3
+ \begin{pmatrix} 0 \\ 1 \end{pmatrix}\mu_4 &= 0
  && (\clubsuit) \\
-x_1 \le 0,\quad x_1 - 2 \le 0,\quad -x_2 \le 0,\quad x_2 - 1 &\le 0
  && (g_1)\text{--}(g_4) \\
\mu_1(-x_1) = 0,\ \mu_2(x_1-2) = 0,\ \mu_3(-x_2) = 0,\ \mu_4(x_2-1) &= 0
  && (\spadesuit) \\
\mu_1 \ge 0,\quad \mu_2 \ge 0,\quad \mu_3 \ge 0,\quad \mu_4 &\ge 0
  && (\heartsuit)
\end{align}

\emph{Provjera gradijenta:} $\nabla f = (2(x_1-3),\, 2(x_2-2)) =
(2x_1-6,\, 2x_2-4)$ \checkmark; gradijenti ograničenja su konstantni vektori
koje čitamo iz koeficijenata: $\nabla g_1 = (-1,0)$, $\nabla g_2 = (1,0)$,
$\nabla g_3 = (0,-1)$, $\nabla g_4 = (0,1)$. \checkmark

\textbf{Tablica svih kombinacija aktivnosti} \str{V3:24}\textbf{.} Zbog
$(\spadesuit)$ svako ograničenje je ili aktivno ($g_i(x) = 0$) ili neaktivno
($g_i(x) < 0$, uz $\mu_i = 0$). S četiri ograničenja to je
$2^4 = 16$ slučajeva:

\begin{center}
\small
\begin{tabular}{ccccc@{\qquad}ccccc}
\toprule
\textbf{sluč.} & $g_1$ & $g_2$ & $g_3$ & $g_4$ &
\textbf{sluč.} & $g_1$ & $g_2$ & $g_3$ & $g_4$ \\
\midrule
1 & 0 & 0 & 0 & 0 & 9  & 1 & 0 & 0 & 0 \\
2 & 0 & 0 & 0 & 1 & 10 & 1 & 0 & 0 & 1 \\
3 & 0 & 0 & 1 & 0 & 11 & 1 & 0 & 1 & 0 \\
4 & 0 & 0 & 1 & 1 & 12 & 1 & 0 & 1 & 1 \\
5 & 0 & 1 & 0 & 0 & 13 & 1 & 1 & 0 & 0 \\
\textbf{6} & \textbf{0} & \textbf{1} & \textbf{0} & \textbf{1}
  & 14 & 1 & 1 & 0 & 1 \\
7 & 0 & 1 & 1 & 0 & 15 & 1 & 1 & 1 & 0 \\
8 & 0 & 1 & 1 & 1 & 16 & 1 & 1 & 1 & 1 \\
\bottomrule
\end{tabular}
\end{center}

\textbf{Legenda za tablicu (sa slajda):} \textbf{1} --- ograničenje je aktivno
($g_i(x) = 0$); \textbf{0} --- ograničenje nije aktivno ($g_i(x) < 0$).

\textbf{Primjer čitanja tablice --- slučaj br.\ 10} (zbog $(\spadesuit)$):
\[
\begin{aligned}
g_1(x) = -x_1 = 0, &\qquad \mu_1 \ge 0 \\
g_2(x) = x_1 - 2 < 0, &\qquad \mu_2 = 0 \\
g_3(x) = -x_2 < 0, &\qquad \mu_3 = 0 \\
g_4(x) = x_2 - 1 = 0, &\qquad \mu_4 \ge 0
\end{aligned}
\]

\textbf{Provjera slučaja br.\ 10} \str{V3:25}\textbf{.} Prema pretpostavci,
imamo $x_1 = 0$, $x_2 = 1$, $\mu_2 = 0$, $\mu_3 = 0$. Gradijentna jednadžba
$(\clubsuit)$ tada postaje
\[
\begin{pmatrix} -6 - \mu_1 \\ -2 - \mu_3 \end{pmatrix} = 0
\]
odnosno imamo $\mu_1 = -6$ i $\mu_3 = -2$. Prema uvjetima $(\heartsuit)$ moramo
imati $\mu_i \ge 0$, pa stoga točka optimuma sigurno \textbf{ne} zadovoljava
pretpostavke iz slučaja br.\ 10.

\emph{Raspis prve komponente:} $2x_1 - 6 - \mu_1 + \mu_2 = 2\cdot0 - 6 - \mu_1
+ 0 = -6 - \mu_1$; izjednačeno s nulom daje $\mu_1 = -6 < 0$ \textbf{--- pad}.

\begin{napomena}[tipfeler na slajdu V3:25]
U ispisu slučaja 10 na slajdu piše $g_3(x) = -x_1 < 0$; treba biti
$g_3(x) = -x_2 < 0$ (usporedi definiciju $g_3$ na str.\ V3:20). Isto tako, u
gradijentnoj jednadžbi druga komponenta bi trebala glasiti $-2 + \mu_4$ (jer je
$\mu_3 = 0$, a $\mu_4$ je taj koji ulazi s $+1$); slajd piše $-2 - \mu_3$.
Zaključak se time ne mijenja --- prva komponenta već sama daje $\mu_1 = -6 < 0$
i ruši slučaj 10.
\end{napomena}

\textbf{Provjera slučaja br.\ 6} \str{V3:26}\textbf{.}
\[
\begin{aligned}
g_1(x) = -x_1 < 0, &\qquad \mu_1 = 0 \\
g_2(x) = x_1 - 2 = 0, &\qquad \mu_2 \ge 0 \\
g_3(x) = -x_2 < 0, &\qquad \mu_3 = 0 \\
g_4(x) = x_2 - 1 = 0, &\qquad \mu_4 \ge 0
\end{aligned}
\]

Prema pretpostavci, imamo $x_1 = 2$, $x_2 = 1$, $\mu_1 = 0$, $\mu_3 = 0$.
Gradijentna jednadžba $(\clubsuit)$ tada postaje
\[
\begin{pmatrix} -2 + \mu_2 \\ -2 + \mu_4 \end{pmatrix} = 0
\]
odnosno imamo $\mu_2 = 2$ i $\mu_4 = 2$.

\emph{Raspis:} prva komponenta je $2x_1 - 6 - \mu_1 + \mu_2 =
2\cdot 2 - 6 - 0 + \mu_2 = -2 + \mu_2$; druga je
$2x_2 - 4 - \mu_3 + \mu_4 = 2\cdot 1 - 4 - 0 + \mu_4 = -2 + \mu_4$.
\checkmark

\textbf{Provjerimo još preostale KKT uvjete:} imamo $g_1(x) = -x_1 = -2 < 0$ i
$g_3(x) = -x_2 = -1 < 0$. \textbf{Svi KKT uvjeti su zadovoljeni.} Zbog jake
dualnosti, ovo ($x_1 = 2$, $x_2 = 1$) je rješenje primarnog problema.

\textbf{Rješenje} \str{V3:27}:
\[
x = \begin{pmatrix} 2 \\ 1 \end{pmatrix}, \qquad
\mu = \begin{pmatrix} 0 \\ 2 \\ 0 \\ 2 \end{pmatrix}
\]

\emph{Optimalna vrijednost:} $p\opt = (2-3)^2 + (1-2)^2 = 1 + 1 = 2$.

\begin{napomena}[kako izbjeći ispitivanje svih 16 slučajeva]
Geometrijski je odmah jasno da će optimum biti u \textbf{gornjem desnom vrhu}
pravokutnika --- jer je ciljna točka $(3,2)$ gore desno. Aktivna će stoga biti
gornja granica $x_1 \le 2$ i gornja granica $x_2 \le 1$, tj.\ slučaj 6.

Tablica sa 16 slučajeva je \emph{sustavan} postupak koji radi i kad geometrija
nije očita; u praksi se počne od najvjerojatnijeg slučaja. Uočite i da svaki
neuspjeli pokušaj (kao slučaj 10) košta samo nekoliko redaka računa.
\end{napomena}
\end{zadatak}

\begin{figure}[H]
\centering
\includegraphics[width=0.80\textwidth]{04-primjer2-kkt.png}
\caption{Vježbe 3, Primjer 2. Dozvoljeni skup je pravokutnik, nivo krivulje su
kružnice oko $(3,2)$, a optimum je vrh $(2,1)$ u kojem su aktivna ograničenja
$g_2$ i $g_4$. Gradijent cilja (zeleno) uravnotežen je zbrojem
$\mu_2\nabla g_2 + \mu_4\nabla g_4$ (crveno).}
\end{figure}

\begin{crtanje}[KKT točka u vrhu pravokutnika]
\textbf{Korak 1.} Nacrtajte osi $x_1$ (od $-1$ do $4{,}5$) i $x_2$ (od $-1$ do
$3{,}5$), jednakim mjerilom.

\textbf{Korak 2 --- prvo ograničenje $g_1$: $-x_1 \le 0$, tj.\ $x_1 \ge 0$.}
Granica je okomita os; dopušteno je \textbf{desno} od nje.

\textbf{Korak 3 --- drugo ograničenje $g_2$: $x_1 \le 2$.} Granica je okomiti
pravac kroz $x_1 = 2$; dopušteno je \textbf{lijevo}. Zajedno s korakom 2 imamo
okomitu traku $0 \le x_1 \le 2$.

\textbf{Korak 4 --- treće i četvrto ograničenje: $0 \le x_2 \le 1$.} Analogno,
vodoravna traka. Presjek dviju traka je \textbf{pravokutnik}
$[0,2]\times[0,1]$ --- osjenčajte ga. To je dozvoljeni skup (politop, usporedi
odjeljak~\ref{sec:poliedar}).

\textbf{Korak 5 --- ciljna točka.} Označite križićem točku $(3,2)$. Ona je
\emph{izvan} pravokutnika --- gore desno.

\textbf{Korak 6 --- nivo krivulje.} $f = (x_1-3)^2+(x_2-2)^2 = c$ su kružnice
oko $(3,2)$ polumjera $\sqrt c$. Nacrtajte ih za $c = 0{,}5;\ 1;\ 2;\ 4$.
Uočite da kružnica $c = 2$ (polumjer $\approx 1{,}41$) prva dodiruje
pravokutnik.

\textbf{Korak 7 --- optimum.} Dodirna točka je vrh $(2,1)$ --- označite je punom
točkom. Provjera: udaljenost od $(3,2)$ do $(2,1)$ je $\sqrt{1+1}=\sqrt2$, pa
je $f = 2$. \checkmark

\textbf{Korak 8 --- gradijent cilja.} Iz $(2,1)$ nacrtajte
$\nabla f = (-2,-2)$: dvije jedinice lijevo, dvije dolje --- dakle
\emph{u} pravokutnik (prema unutra).

\textbf{Korak 9 --- gradijenti aktivnih ograničenja.} Iz iste točke nacrtajte
$\nabla g_2 = (1,0)$ (udesno) i $\nabla g_4 = (0,1)$ (gore). Skalirajte ih s
$\mu_2 = 2$ odn.\ $\mu_4 = 2$, dakle po dvije jedinice.

\textbf{Korak 10 --- provjera KKT-a okom.} Zbroj crvenih strelica je
$(2,0) + (0,2) = (2,2)$, što je točno $-\nabla f$. Uvjet $(\clubsuit)$ kaže
upravo to: $\nabla f + \mu_2\nabla g_2 + \mu_4\nabla g_4 = 0$.

\textbf{Što graf znači.} U vrhu politopa aktivna su dva ograničenja, pa
$-\nabla f$ mora ležati \emph{unutar} kuta razapetog njihovim gradijentima ---
to je stožac iz odjeljka o geometrijskoj interpretaciji \str{3:56}. Da je ciljna
točka bila npr.\ $(3;\ 0{,}5)$, aktivno bi bilo samo $g_2$, i optimum bi bio na
stranici, a ne u vrhu.
\end{crtanje}

% =====================================================================
\section{Numeričko rješavanje i dualni problem}
\label{sec:primjer2-dual}

Ovaj odjeljak nastavlja Zadatak~\ref{zad:primjer2} podzadacima c) i d).

\subsection{c) Rješavanje u MATLAB-u}
\str{V3:28--33}

\textbf{Prevođenje u QP standardnu formu} \str{V3:28}\textbf{.} MATLAB-ova
funkcija \texttt{quadprog} očekuje problem oblika
$\min \tfrac12 x\T Q x + c\T x$ uz $Gx \le h$. Zato ciljnu funkciju treba
raspisati:

\[
\text{minimiziraj}_x \quad
\frac{1}{2}\begin{pmatrix}x_1\\x_2\end{pmatrix}\T
\underbrace{\begin{pmatrix}2&0\\0&2\end{pmatrix}}_{Q}
\begin{pmatrix}x_1\\x_2\end{pmatrix}
+ \underbrace{\begin{pmatrix}-6 & -4\end{pmatrix}}_{c\T}
\begin{pmatrix}x_1\\x_2\end{pmatrix} + 13
\]
\[
\text{uz ogr.}\quad
\underbrace{\begin{pmatrix}
-1 & 0 \\ 1 & 0 \\ 0 & -1 \\ 0 & 1
\end{pmatrix}}_{G}
\begin{pmatrix}x_1\\x_2\end{pmatrix}
\le
\underbrace{\begin{pmatrix}0\\2\\0\\1\end{pmatrix}}_{h}
\]

\emph{Provjera raspisa ciljne funkcije:}
\begin{align}
(x_1-3)^2 + (x_2-2)^2
&= x_1^2 - 6x_1 + 9 + x_2^2 - 4x_2 + 4 \\
&= \underbrace{x_1^2 + x_2^2}_{=\;\frac12(2x_1^2+2x_2^2)\;=\;\frac12 x\T Qx}
 \;\underbrace{-\,6x_1 - 4x_2}_{=\;c\T x}\; + \;13
\end{align}
\checkmark \; Konstanta $13$ ne utječe na $\argmin$, pa je \texttt{quadprog} ne
prima --- ali je treba \textbf{dodati natrag} na dobivenu optimalnu vrijednost.
(To je odgovor na pitanje „zašto $+13$?” s komentara u kodu.)

\textbf{Kod sa slajda} \str{V3:30}:

\begin{verbatim}
Q=[2, 0;  0, 2];
c=[-6; -4];
G=[-1,0;  1,0;  0,-1;  0,1];
h=[0;2;0;1];

[x_opt,f_opt,EXITFLAG,OUTPUT,mu_opt] = quadprog(Q,c,G,h);

x_opt              % optimalne primarne varijable
mu_opt.ineqlin     % optimalne dualne varijable

f_opt+13           % vrijednost funkcije cilja u optimiumu (zasto +13 ?)
\end{verbatim}

\textbf{Ispis rezultata} \str{V3:31}:

\begin{verbatim}
>> x_opt
x_opt =
    2.0000
    1.0000

>> mu_opt.ineqlin
ans =
         0
    2.0000
         0
    2.0000

>> f_opt+13
ans =
    2.0000
\end{verbatim}

\textbf{Poklapanje s ručnim rješenjem:} $x\opt = (2,1)$,
$\mu\opt = (0,2,0,2)$, $p\opt = 2$ --- točno kao u podzadatku b). \checkmark

\textbf{Isti problem preko YALMIP-a} \str{V3:32}\textbf{.} YALMIP je
„modelirajući sloj”: problem se piše prirodnim zapisom, a on ga sam prevede
odabranom rješavaču (ovdje opet \texttt{quadprog}).

\begin{verbatim}
yalmip('clear')
x=sdpvar(2,1);

Con=[0<=x(1)<=2, 0<=x(2)<=1];

Objective=(x(1)-3)^2+(x(2)-2)^2;

% neke opcije YALMIP i "solver"
options = sdpsettings('verbose',1,'solver','quadprog',...
                      'quadprog.maxiter',100);

% Poziv za rjesavanje
sol = optimize(Con,Objective,options);

% Analiza "error zastavica"
if sol.problem == 0
 % Definiraj i ispintaj rjesenje
 x_opt = value(x)     % optimalne primarne
 mu_opt=dual(Con)     % optimalne dualne varijable
else
 display('Nesto je krivo!');
 sol.info
 yalmiperror(sol.problem)
end
\end{verbatim}

\textbf{Ispis} \str{V3:33}\textbf{:} \texttt{x\_opt = [2.0000; 1.0000]},
\texttt{mu\_opt = [0; 2.0000; 0; 2.0000]} --- identično. \checkmark

\begin{napomena}[quadprog vs.\ YALMIP]
\texttt{quadprog} traži da problem \emph{ručno} prevedete u matrice
$Q, c, G, h$ --- što je izvor grešaka, ali daje potpunu kontrolu. YALMIP
prihvaća zapis blizak matematičkom (\texttt{0<=x(1)<=2}) i sam radi prijevod;
zato je pogodniji za veće modele. Obratite pažnju da YALMIP dualne varijable
vraća funkcijom \texttt{dual(Con)}, a \texttt{quadprog} u polju
\texttt{mu\_opt.ineqlin}.
\end{napomena}

\subsection{d) Postavljanje i rješavanje dualnog problema}
\str{V3:34--39}

\textbf{Primarni problem} (u QP formi, bez konstante):
\[
\min_x \; \tfrac12 x\T Qx + c\T x \qquad\text{uz ogr.}\quad Gx - h \le 0
\]

\textbf{Uočimo} (sa slajda): $Q \succ 0$, pa stoga i $Q^{-1}$ postoji ($Q$ nema
svojstvenu vrijednost jednaku nuli jer su sve striktno veće od nule).

\textbf{Korak 1 --- Lagrangian:}
\[
L(x,\mu) = \tfrac12 x\T Qx + c\T x + \mu\T(Gx - h)
         = \tfrac12 x\T Qx + (c\T + \mu\T G)x - \mu\T h
\]

\textbf{Korak 2 --- minimizacija po $x$} \str{V3:35}:
\[
\ell(\mu) = \inf_x L(x,\mu) \;\to\;
\nabla_x L(x,\mu) = 0 \;\to\;
Q\tilde x + c + G\T\mu = 0 \;\to\;
\tilde x = -Q^{-1}(c + G\T\mu)
\]

\emph{Raspis gradijenta:} za kvadratnu formu vrijedi
$\nabla_x(\tfrac12 x\T Qx) = Qx$ (uz simetričan $Q$), a
$\nabla_x\big((c\T+\mu\T G)x\big) = c + G\T\mu$. Član $-\mu\T h$ ne ovisi o
$x$. Izjednačavanjem s nulom i množenjem s $Q^{-1}$ slijeva dobivamo $\tilde x$.
\checkmark

\textbf{Korak 3 --- uvrštavanje natrag} \str{V3:36}:
\begin{align}
\ell(\mu) &= \tfrac12\tilde x\T Q\tilde x + (c\T + \mu\T G)\tilde x - \mu\T h \\
&= -\tfrac12\mu\T GQ^{-1}G\T\mu
   + \big(-c\T Q^{-1}G\T - h\T\big)\mu
   - \tfrac12 c\T Q^{-1}c
\end{align}

\emph{Raspis (na slajdu nije prikazan).} Označimo $z := c + G\T\mu$, pa je
$\tilde x = -Q^{-1}z$ i $(c\T+\mu\T G) = z\T$:
\begin{align}
\tfrac12\tilde x\T Q\tilde x
  &= \tfrac12 z\T Q^{-1}QQ^{-1}z = \tfrac12 z\T Q^{-1}z \\
z\T\tilde x &= -z\T Q^{-1}z \\
\Rightarrow\quad \ell(\mu)
  &= \tfrac12 z\T Q^{-1}z - z\T Q^{-1}z - \mu\T h
   = -\tfrac12 z\T Q^{-1}z - \mu\T h
\end{align}
Sada raspišemo $z\T Q^{-1}z$:
\[
(c + G\T\mu)\T Q^{-1}(c + G\T\mu)
= c\T Q^{-1}c + 2c\T Q^{-1}G\T\mu + \mu\T GQ^{-1}G\T\mu
\]
Uvrštavanjem:
\[
\ell(\mu) = -\tfrac12 c\T Q^{-1}c - c\T Q^{-1}G\T\mu
            - \tfrac12\mu\T GQ^{-1}G\T\mu - h\T\mu
\]
što je točno izraz sa slajda nakon grupiranja linearnih članova. \checkmark

\textbf{Dualni problem} \str{V3:37}:
\[
\max_{\mu \ge 0}\;
-\tfrac12\mu\T GQ^{-1}G\T\mu + \big(-c\T Q^{-1}G\T - h\T\big)\mu
- \tfrac12 c\T Q^{-1}c
\]

\textbf{Tvrdnja sa slajda:} budući da je $-GQ^{-1}G\T \prec 0$, dualni problem
je problem maksimizacije konkavne funkcije s konveksnim ograničenjima
$\to$ \textbf{konveksan optimizacijski problem}.

\begin{napomena}[preciznije o toj tvrdnji]
Općenito vrijedi $-GQ^{-1}G\T \preceq 0$ (negativno \emph{semi}definitno), jer
je $\mu\T GQ^{-1}G\T\mu = (G\T\mu)\T Q^{-1}(G\T\mu) \ge 0$ uz $Q^{-1} \succ 0$
--- ali je nula kad god je $G\T\mu = 0$.

U ovom konkretnom primjeru $G$ ima $4$ retka, a samo $2$ stupca, pa $G\T$ ima
netrivijalnu jezgru i matrica $GQ^{-1}G\T$ \emph{jest} singularna. Konkretno,
uz $Q^{-1} = \tfrac12 I$:
\[
GQ^{-1}G\T = \tfrac12 GG\T = \tfrac12
\begin{pmatrix}
 1 & -1 & 0 & 0 \\
-1 & 1 & 0 & 0 \\
 0 & 0 & 1 & -1 \\
 0 & 0 & -1 & 1
\end{pmatrix}
\]
čije su svojstvene vrijednosti $0, 0, 1, 1$. Dakle striktna negativna
definitnost sa slajda je preoštra; ispravno je $\preceq 0$. \textbf{Zaključak
ostaje isti} --- za konkavnost je dovoljna semidefinitnost, pa je dualni
problem i dalje konveksan.
\end{napomena}

\textbf{Kod za rješavanje dualnog problema} \str{V3:38}:

\begin{verbatim}
Q=[2, 0;  0, 2];
c=[-6; -4];
G=[-1,0;  1,0;  0,-1;  0,1];
h=[0;2;0;1];

Q_dual=G*inv(Q)*G';
c_dual=G*inv(Q)*c+h;
G_dual=-eye(4);
h_dual=zeros(4,1);

[mu_opt_Dual,f_opt_Dual] = quadprog(Q_dual,c_dual,G_dual,h_dual);

mu_opt_Dual   % optimalne dualne varijable

x_opt_Dual=-inv(Q)*(c+G'*mu_opt_Dual)   % optimalne primarne varijable

% vrijednost dualne funkcije cilja u optimiumu:
f_opt_Dual=(-f_opt_Dual+13)-0.5*c'*inv(Q)*c
\end{verbatim}

\textbf{Kako je kod dobiven iz formule} --- vrijedi ga pročitati redak po redak:

\begin{itemize}
  \item \texttt{quadprog} \emph{minimizira}, a mi \emph{maksimiziramo}, pa se
        koristi trik $\max \ell \Leftrightarrow \min(-\ell)$
        (odjeljak~\ref{sec:trikovi}). Zato \texttt{Q\_dual} $= GQ^{-1}G\T$ (bez
        minusa) i \texttt{c\_dual} $= GQ^{-1}c + h$.
  \item Ograničenje $\mu \ge 0$ piše se u obliku $G_{\text{dual}}\mu \le
        h_{\text{dual}}$ uz $G_{\text{dual}} = -I$ i $h_{\text{dual}} = 0$, tj.\
        $-\mu \le 0$.
  \item \texttt{x\_opt\_Dual} rekonstruira primarnu varijablu iz formule
        $\tilde x = -Q^{-1}(c + G\T\mu)$ --- \textbf{primarno rješenje se dobiva
        iz dualnog}, što je moguće upravo zbog jake dualnosti.
  \item Zadnji redak vraća predznak i dodaje konstante: $+13$ (izbačena iz
        cilja) i $-\tfrac12 c\T Q^{-1}c$ (član koji \texttt{quadprog} ne vidi).
\end{itemize}

\textbf{Ispis} \str{V3:39}:

\begin{verbatim}
mu_opt_Dual =        x_opt_Dual =        f_opt_Dual =
    0.0000               2.0000              2.0000
    2.0000               1.0000
    0.0000
    2.0000
\end{verbatim}

\textbf{Provjera jake dualnosti:} $d\opt = 2 = p\opt$. \checkmark Dualne
varijable dobivene rješavanjem \emph{dualnog} problema identične su onima koje
je \texttt{quadprog} vratio kao nusprodukt rješavanja \emph{primarnog} problema.

% =====================================================================
\section{Dualna dekompozicija: distribuirano rješavanje}
\label{sec:dekompozicija}
\str{24--27}

\subsection{Strukturirani problem}
\str{25}

\begin{equation}
\label{eq:strukturiran}
\begin{aligned}
\min_{x_1,x_2} \quad & f_1(x_1) + f_2(x_2) \\
\text{uz ogr.} \quad & x_1 \in \Feas_1, \; x_2 \in \Feas_2 \\
& h_1(x_1) + h_2(x_2) = 0 && (\spadesuit)
\end{aligned}
\end{equation}

\textbf{Opažanje sa slajda:} kad ne bi bilo ograničenja $(\spadesuit)$ koje
„povezuje” $x_1$ i $x_2$, ovaj optimizacijski problem bi se sveo na
\textbf{dva neovisna} optimizacijska problema.

\textbf{Struktura problema.} Funkcija cilja je \emph{separabilna} (zbroj
članova, svaki ovisi o svojoj varijabli), lokalna ograničenja su odvojena, i
samo \emph{jedno} ograničenje $(\spadesuit)$ spaja dva dijela. Takva se
struktura pojavljuje svugdje gdje više jedinica dijeli zajednički resurs.

\subsection{Dualni problem}
\str{26}

\begin{align}
\max_\lambda\; \ell(\lambda)
&= \inf_{x_1\in\Feas_1,\,x_2\in\Feas_2}
   f_1(x_1) + f_2(x_2) + \lambda\T\big(h_1(x_1) + h_2(x_2)\big) \\
&= \inf_{x_1\in\Feas_1}\Big(f_1(x_1) + \lambda\T h_1(x_1)\Big)
 + \inf_{x_2\in\Feas_2}\Big(f_2(x_2) + \lambda\T h_2(x_2)\Big)
\end{align}

\emph{Zašto se infimum smije rastaviti:} nakon uvrštavanja $\lambda$ izraz se
raspada na dva dijela, od kojih prvi ovisi \emph{samo} o $x_1$, a drugi
\emph{samo} o $x_2$; varijable više nisu spojene nikakvim zajedničkim
ograničenjem. Minimum zbroja neovisnih članova je zbroj minimuma.

\textbf{Ključno opažanje sa slajda:} \emph{Uočite da smo dualizirali samo
globalno ograničenje $(\spadesuit)$ dok su lokalna ograničenja
$x_1 \in \Feas_1$, $x_2 \in \Feas_2$ ostavljena u definiciji od $\ell(\lambda)$.}

\begin{napomena}[djelomična dualizacija]
Nije nužno dualizirati \emph{sva} ograničenja. Bira se ono čije uklanjanje
razbija strukturu problema --- ovdje globalno ograničenje. Lokalna ostaju kao
ograničenja unutarnjih podproblema. To je bit tehnike.
\end{napomena}

\subsection{Algoritam}
\str{27}

Za zadani fiksni $\lambda$ imamo \textbf{dva neovisna podproblema} (mogu se
rješavati odvojeno / paralelno):

\begin{itemize}
  \item \textbf{podproblem 1:}
        $\inf_{x_1\in\Feas_1}\big(f_1(x_1) + \lambda\T h_1(x_1)\big)$,
        koji ima opt.\ vrijednost $\ell_1(\lambda)$;
  \item \textbf{podproblem 2:}
        $\inf_{x_2\in\Feas_2}\big(f_2(x_2) + \lambda\T h_2(x_2)\big)$,
        koji ima opt.\ vrijednost $\ell_2(\lambda)$.
\end{itemize}

\textbf{Dualni problem} (koji rješava i primarni, uz jaku dualnost):
\[
\max_\lambda\; \ell(\lambda) = \ell_1(\lambda) + \ell_2(\lambda)
\]

\textbf{Rješenje dualnog problema se postiže kad je $h_1(x_1) + h_2(x_2) = 0$.}

\begin{napomena}[kako to izgleda kao algoritam]
Postupak je iterativan i ima strukturu „koordinator + jedinice”:
\begin{enumerate}
  \item koordinator objavi trenutni $\lambda$;
  \item svaka jedinica \emph{neovisno} riješi svoj mali problem i javi svoj
        $h_i(x_i)$;
  \item koordinator pogleda „neravnotežu” $h_1 + h_2$ i prilagodi $\lambda$
        (poveća ga ako je zbroj pozitivan, smanji ako je negativan);
  \item ponovi dok neravnoteža ne padne na nulu.
\end{enumerate}
Uočite koliko se malo informacija razmjenjuje: jedinice \textbf{ne moraju
otkriti} svoje funkcije $f_i$, $h_i$ ni svoje skupove $\Feas_i$ --- samo jedan
broj. To je ključno u primjeni koja slijedi.
\end{napomena}

% =====================================================================
\section{Vježbe 3, Primjer 1: tržište i formiranje cijena}
\label{sec:trziste}
\str{V3:3--18}

Ovo je izravna primjena dualne dekompozicije --- i objašnjenje najave iz
poglavlja~\ref{ch:uvod} \str{1:34}: \emph{matematička formulacija tržišta i
kreiranja tržišnih cijena je Lagrangeova dualnost}.

\subsection{Postavka: maksimizacija društvene dobrobiti}
\str{V3:4}

\textbf{Oznake:}
\begin{itemize}
  \item Troškovi proizvodnje za $i$-tog proizvođača za proizvodnju $p_i$ [MWh]
        energije: $C_i(p_i)$.
  \item Funkcija benefita za $j$-tog potrošača kod potrošnje $d_j$ [MWh]
        energije: $B_j(d_j)$.
\end{itemize}

\textbf{Maksimizacija društvene dobrobiti:}
\begin{equation}
\label{eq:trziste}
\begin{aligned}
\min_{p_1,\dots,p_n,\,d_1,\dots,d_m} \quad
& \sum_{i=1}^{n}C_i(p_i) - \sum_{j=1}^{m}B_j(d_j)
  && (=\ \max\ \text{društvena dobrobit}) \\
\text{uz ogr.} \quad
& p_i \in P_i, \; i = 1,\dots,n
  && \text{(lokalna ograničenja proizvođača)} \\
& d_j \in D_j, \; j = 1,\dots,m
  && \text{(lokalna ograničenja potrošača)} \\
& \sum_{i=1}^{n}p_i = \sum_{j=1}^{m}d_j
  && \text{(balans potrošnje i proizvodnje)}
\end{aligned}
\end{equation}

\textbf{Primjeri lokalnih ograničenja} (sa slajda):
\[
P_i := \{\, p \mid \underline{p_i} \le p \le \overline{p_i} \,\}, \qquad
D_j := \{\, d \mid \underline{d_j} \le d \le \overline{d_j} \,\}
\]
--- dakle svaka elektrana ima minimalnu i maksimalnu snagu, svaki potrošač
minimalnu i maksimalnu potrošnju.

\textbf{Kako čitati funkciju cilja.} \emph{Društvena dobrobit} je ukupna korist
potrošača umanjena za ukupni trošak proizvodnje,
$\sum_j B_j(d_j) - \sum_i C_i(p_i)$. Nju želimo \emph{maksimizirati}, pa
minimiziramo njezinu negaciju --- odatle poredak članova u
\eqref{eq:trziste} (trik $\max f \Leftrightarrow \min(-f)$,
odjeljak~\ref{sec:trikovi}).

\textbf{Pretpostavka konveksnosti} \str{V3:6}\textbf{:} $C_i(\cdot)$ konveksne
funkcije, $B_j(\cdot)$ \textbf{konkavne} funkcije, $P_i$, $D_j$ konveksni
skupovi.

\begin{napomena}[zašto baš takve pretpostavke]
\begin{itemize}
  \item $C_i$ konveksna: trošak proizvodnje raste \emph{sve brže} --- svaki
        sljedeći MWh je skuplji (uključuju se sve skuplji agregati). To je
        realistično.
  \item $B_j$ konkavna: korist raste \emph{sve sporije} --- zakon opadajuće
        granične korisnosti. Prvi MWh vrijedi mnogo, deseti mnogo manje.
  \item Uz to je $-B_j$ konveksna, pa je \textbf{cijela funkcija cilja
        konveksna} --- problem \eqref{eq:trziste} je konveksan i vrijedi jaka
        dualnost.
\end{itemize}
\end{napomena}

\subsection{Dualni problem}
\str{V3:5--9}

\textbf{Primarni problem} \str{V3:5}:
\[
\min_{p_i\in P_i,\,d_j\in D_j}\;
\sum_{i=1}^{n}C_i(p_i) - \sum_{j=1}^{m}B_j(d_j)
\qquad\text{uz ogr.}\quad \sum_{i=1}^{n}p_i = \sum_{j=1}^{m}d_j
\]

\textbf{Dualni problem} \str{V3:6}:
\[
\max_{\lambda\in\R}\; \ell(\lambda), \quad\text{gdje je}\quad
\ell(\lambda) = \min_{p_i\in P_i,\,d_j\in D_j}
\sum_{i=1}^{n}C_i(p_i) - \sum_{j=1}^{m}B_j(d_j)
+ \lambda\left(\sum_{j=1}^{m}d_j - \sum_{i=1}^{n}p_i\right)
\]

\textbf{Napomena sa slajda} \str{V3:7}\textbf{:} \emph{Uočite da smo dualizirali
samo globalno ograničenje $\sum_j d_j - \sum_i p_i = 0$ dok su ostala lokalna
ograničenja ostavljena u definiciji od $\ell(\lambda)$.}

\textbf{Opažanje 1} \str{V3:8}\textbf{:} Lagrangeova dualna funkcija
$\ell(\lambda)$ se za fiksan $\lambda$ može rastaviti u \textbf{$n + m$
odvojenih minimizacijskih problema}.

\textbf{Opažanje 2} \str{V3:9}\textbf{:} $\max_{\lambda\in\R}\ell(\lambda)$ se
postiže kad je $\sum_{j=1}^{m}d_j = \sum_{i=1}^{n}p_i$ ((sub)gradijent od
$\ell(\lambda)$ je nula).

\subsection{Rastav na lokalne probleme}
\str{V3:10}

\begin{center}
\begin{tabular}{P{6.4cm}P{6.4cm}}
\toprule
\textbf{Lokalna minimizacija proizvođača} &
\textbf{Lokalna minimizacija potrošača} \\
$\to$ maksimizacija svoje zarade &
$\to$ maksimizacija svog benefita \\
\midrule
$\displaystyle\min_{p_1\in P_1} C_1(p_1) - \lambda p_1$ &
$\displaystyle\min_{d_1\in D_1} \lambda d_1 - B_1(d_1)$ \\[6pt]
$\displaystyle\min_{p_2\in P_2} C_2(p_2) - \lambda p_2$ &
$\displaystyle\min_{d_2\in D_2} \lambda d_2 - B_2(d_2)$ \\[2pt]
\quad$\vdots$ & \quad$\vdots$ \\[2pt]
$\displaystyle\min_{p_n\in P_n} C_n(p_n) - \lambda p_n$ &
$\displaystyle\min_{d_m\in D_m} \lambda d_m - B_m(d_m)$ \\
\bottomrule
\end{tabular}
\end{center}

\textbf{Interpretacija --- ovo je srce cijelog primjera:}

\begin{itemize}
  \item $C_i(p_i) - \lambda p_i$ je \textbf{negativna zarada} proizvođača:
        $\lambda p_i$ je prihod od prodaje $p_i$ MWh po cijeni $\lambda$, a
        $C_i(p_i)$ je trošak. Minimizacija tog izraza = maksimizacija profita.
  \item $\lambda d_j - B_j(d_j)$ je \textbf{negativna korist} potrošača:
        $\lambda d_j$ je ono što plaća, $B_j(d_j)$ je korist koju ima.
  \item Dakle: \textbf{svaki sudionik gleda samo svoj interes}, koristeći samo
        \emph{svoju} funkciju i \emph{jedan zajednički broj} $\lambda$.
\end{itemize}

\begin{napomena}[tipfeler na slajdovima V3:10--11]
U stupcu potrošača na slajdu svugdje piše $B_1(\cdot)$, i za drugog i za
$m$-tog potrošača. Ispravno je $B_2(d_2)$, \dots, $B_m(d_m)$ --- svaki potrošač
ima svoju funkciju benefita.
\end{napomena}

\subsection{Uloga operatora tržišta}
\str{V3:11}

\begin{center}
\begin{tabular}{P{13.0cm}}
\toprule
\textbf{Operator tržišta:} \quad
$\displaystyle\max_{\lambda\in\R}\ell(\lambda)
\quad\Longleftrightarrow\quad
\text{odredi }\lambda:\; \sum_{j=1}^{m}d_j\opt = \sum_{i=1}^{n}p_i\opt$ \\
\bottomrule
\end{tabular}
\end{center}

\textbf{Racionalno ponašanje sudionika} (max svoj benefit):
\[
p_i\opt = \argmin_{p_i\in P_i} C_i(p_i) - \lambda p_i, \qquad
d_j\opt = \argmin_{d_j\in D_j} \lambda d_j - B_j(d_j)
\]

\begin{teorem}[Glavni rezultat \str{V3:11}]
$\lambda\opt$ koji rješava ovaj problem $=$ \textbf{tržišna cijena energije}.
\end{teorem}

\textbf{Zašto je to fascinantno.} Dualna varijabla, uvedena kao čisto
matematička pomoć, ispada \textbf{cijena}. Nitko ne mora znati tuđe troškove ni
funkcije koristi; svatko optimira svoj interes uz objavljenu cijenu, a operator
samo podešava cijenu dok se ponuda i potražnja ne izjednače. Rezultat je
\textbf{globalno optimalna} raspodjela --- jer je to rješenje dualnog problema,
a uz jaku dualnost i primarnog.

Ovo je ujedno i matematički sadržaj metafore o „nevidljivoj ruci tržišta”.

\subsection{Marginalni troškovi i tržišna cijena}
\str{V3:12}

\textbf{Tvrdnja.} Neka je $\lambda$ takav da je $p_i\opt \in$ unutrašnjost od
$P_i$, $d_j\opt \in$ unutrašnjost od $D_j$; tada imamo
\[
\frac{\dd C_i(p_i\opt)}{\dd p_i} = \lambda, \qquad
\frac{\dd B_j(d_j\opt)}{\dd d_j} = \lambda
\]

\emph{Odakle:} ako je optimum u \emph{unutrašnjosti} lokalnog skupa, lokalna
ograničenja su neaktivna, pa vrijedi obični uvjet stacionarnosti
(odjeljak~\ref{sec:uvjeti-bez-ogr}): derivacija cilja po varijabli je nula.
Za proizvođača: $\dfrac{\dd}{\dd p_i}\big(C_i(p_i) - \lambda p_i\big)
= C_i'(p_i) - \lambda = 0$. \checkmark Analogno za potrošača:
$\lambda - B_j'(d_j) = 0$.

\textbf{Zaključak sa slajda:} \emph{društvena dobrobit je maksimizirana kad svi
potrošači i proizvođači postave nivo proizvodnje/potrošnje tako da su
inkrementalni (marginalni) troškovi/benefiti jednaki tržišnoj cijeni $\lambda$.}

\begin{napomena}[marginalni trošak]
\textbf{Marginalni (inkrementalni) trošak} $\dd C/\dd p$ je trošak proizvodnje
\emph{jednog dodatnog} MWh. Uvjet kaže: proizvodi sve dok te sljedeći MWh košta
manje nego što za njega dobivaš; stani kad se izjednače. Isto vrijedi za
potrošača s koristi. \textbf{U optimumu su svi marginalni troškovi i sve
marginalne koristi jednaki --- i to baš cijeni $\lambda$.}
\end{napomena}

\subsection{Krivulje ponude i potražnje}
\str{V3:13--14}

Iz uvjeta marginalnosti dobivaju se \textbf{krivulje ponude i potražnje}:
\[
\frac{\dd C_i(p_i)}{\dd p_i} = \lambda
\;\Leftrightarrow\; p_i = \gamma_i^p(\lambda)
\;\Leftrightarrow\; \lambda = \beta_i^p(p_i)
\]
\[
\frac{\dd B_j(d_j)}{\dd d_j} = \lambda
\;\Leftrightarrow\; d_j = \gamma_j^d(\lambda)
\;\Leftrightarrow\; \lambda = \beta_j^d(d_j)
\]

To je ista relacija napisana na tri načina: kao jednadžba, kao „koliko
proizvodim pri cijeni $\lambda$” ($\gamma$), i kao „pri kojoj cijeni proizvodim
$p_i$” ($\beta$) --- dakle inverzne funkcije jedna drugoj.

\textbf{Određivanje tržišne cijene u praksi} \str{V3:14}\textbf{.} Cijena
$\lambda\opt$ nalazi se na sjecištu \textbf{agregirane krivulje ponude}
$\tilde\gamma^p(\lambda) := \sum_i\gamma_i^p(\lambda)$ s \textbf{agregiranom
krivuljom potražnje} $\tilde\gamma^d(\lambda) := \sum_j\gamma_j^d(\lambda)$:
\[
\sum_{i=1}^{n}p_i\opt = \sum_{i=1}^{n}\gamma_i^p(\lambda\opt)
= \tilde\gamma^p(\lambda\opt) = \tilde\gamma^d(\lambda\opt)
= \sum_{j=1}^{m}\gamma_j^d(\lambda\opt) = \sum_{j=1}^{m}d_j\opt
\]

\textbf{Napomena sa slajda:} generalizacija na slučajeve kad ne vrijedi da je
$p_i\opt \in$ unutrašnjost od $P_i$, $d_j\opt \in$ unutrašnjost od $D_j$ je
moguća. Lokalna ograničenja mogu se uključiti u krivulje ponude/potražnje.

\textbf{Slika sa slajda} \str{V3:15}\textbf{.} Prikazana su tri grafa:

\begin{itemize}
  \item gore lijevo „\emph{Supply}”: os $\lambda$ okomito, os $p$ vodoravno;
        nekoliko stepenastih rastućih krivulja (po jedna za svakog
        proizvođača);
  \item gore desno „\emph{Demand}”: os $\lambda$ okomito, os $d$ vodoravno;
        nekoliko stepenastih \emph{padajućih} krivulja;
  \item dolje: dvije strelice vode prema trećem grafu, na kojem su
        \textbf{agregirane} krivulje --- rastuća (ponuda) i padajuća
        (potražnja). Njihovo sjecište označeno je crvenom točkom; na osi
        $\lambda$ očitava se $\lambda^{*}$, a na vodoravnoj osi („\emph{power}”)
        očitava se $\sum p_i^{*} = \sum d_j^{*}$.
\end{itemize}

\begin{crtanje}[određivanje tržišne cijene]
\textbf{Korak 1.} Nacrtajte koordinatni sustav: vodoravno \emph{količina
energije} (snaga), okomito \emph{cijena} $\lambda$.

\textbf{Korak 2 --- krivulja ponude.} Poredajte proizvođače po marginalnom
trošku, od najjeftinijeg prema najskupljem. Crtajte stepenice slijeva nadesno:
širina stepenice je kapacitet tog proizvođača, visina njegov marginalni trošak.
Dobivate \textbf{rastuću} stepenastu krivulju --- to je agregirana ponuda
$\tilde\gamma^p$.

\emph{Zašto raste:} prvo se uključuju najjeftiniji izvori; što više energije
treba, to skuplje agregate moramo pokrenuti.

\textbf{Korak 3 --- krivulja potražnje.} Isto, ali za potrošače poredane po
marginalnoj koristi, od najviše prema najnižoj. Dobivate \textbf{padajuću}
stepenastu krivulju $\tilde\gamma^d$.

\emph{Zašto pada:} najprije se zadovoljavaju potrebe koje potrošač najviše
cijeni; dodatna potrošnja vrijedi mu sve manje.

\textbf{Korak 4 --- sjecište.} Dvije se krivulje sijeku u jednoj točki.
Označite je. Iz nje povucite vodoravnu isprekidanu crtu do osi $\lambda$ ---
tamo očitajte \textbf{tržišnu cijenu} $\lambda^{*}$. Povucite i okomitu
isprekidanu crtu do vodoravne osi --- tamo očitajte \textbf{količinu}
$\sum p_i^{*} = \sum d_j^{*}$.

\textbf{Korak 5 --- što graf znači.} Lijevo od sjecišta ponuda je jeftinija od
onoga što potrošači vrednuju --- svaka takva transakcija povećava društvenu
dobrobit i \emph{treba} se dogoditi. Desno od sjecišta proizvodnja bi koštala
više nego što vrijedi --- takve transakcije smanjuju dobrobit. Sjecište je
točka gdje se to izjednači, dakle upravo optimum problema
\eqref{eq:trziste}, a $\lambda^{*}$ je dualna varijabla balansnog ograničenja.
\end{crtanje}

\begin{figure}[H]
\centering
\includegraphics[width=0.74\textwidth]{04-ponuda-potraznja.png}
\caption{Tržišna cijena $\lambda^\star$ kao sjecište agregirane krivulje ponude
(rastuća) i potražnje (padajuća). Rekonstrukcija sheme s Vježbi 3, str.\ 15.}
\end{figure}

\textbf{Stvarni podaci sa slajdova} \str{V3:17--18}\textbf{.} Slajdovi
prikazuju agregirane krivulje burze \textbf{APX} (Amsterdam Power Exchange) za
dva termina istog dana:

\begin{itemize}
  \item \textbf{30.\ siječnja 2015., 2 h ujutro} --- noćni sat, niska
        potražnja: krivulja potražnje je pomaknuta ulijevo, pa je sjecište
        \emph{niže} na osi cijena;
  \item \textbf{30.\ siječnja 2015., 19 h} --- večernji vršni sat, visoka
        potražnja: krivulja potražnje je pomaknuta udesno, sjecište je na
        \emph{strmom} dijelu krivulje ponude, pa je cijena znatno viša.
\end{itemize}

\textbf{Poruka:} teorija s prethodnih slajdova doslovno je mehanizam po kojem
europske burze električne energije svakodnevno formiraju cijenu.

% =====================================================================
\section{Analiza osjetljivosti}
\label{sec:osjetljivost}
\str{28--34}

Zadnja tema poglavlja odgovara na pitanje: \emph{koliko se optimum promijeni ako
malo pomaknemo ograničenja?}

\subsection{Perturbiran problem}
\str{29--31}

\textbf{Polazni par problema} \str{29}:

\begin{center}
\begin{tabular}{P{6.3cm}P{7.0cm}}
\toprule
\textbf{Primarni problem} & \textbf{Dualni problem} \\
\midrule
$\displaystyle\min_x f(x)$ &
$\displaystyle\max_{\lambda,\mu}\;
\ell(\lambda,\mu) := \inf_x\big(f(x) + \lambda\T h(x) + \mu\T g(x)\big)$ \\
uz ogr.\ $h_i(x) = 0$, $i=1,\dots,m$ & uz ogr.\ $\mu \ge 0$ \\
\phantom{uz ogr.\ }$g_j(x) \le 0$, $j=1,\dots,p$ & \\
\bottomrule
\end{tabular}
\end{center}

\textbf{Perturbiran optimizacijski problem} \str{30}:

\begin{center}
\begin{tabular}{P{6.3cm}P{7.0cm}}
\toprule
\textbf{Primarni problem} & \textbf{Dualni problem} \\
\midrule
$\displaystyle\min_x f(x)$ &
$\displaystyle\max_{\lambda,\mu}\;
\ell(\lambda,\mu) - \lambda\T u - \mu\T v$ \\
uz ogr.\ $h_i(x) = u_i$, $i=1,\dots,m$ & uz ogr.\ $\mu \ge 0$ \\
\phantom{uz ogr.\ }$g_j(x) \le v_j$, $j=1,\dots,p$ & \\
\bottomrule
\end{tabular}
\end{center}

gdje su $u = [u_1 \dots u_m]\T$, $v = [v_1 \dots v_p]\T$
\textbf{perturbacijski parametri}.

\emph{Odakle dualni oblik:} Lagrangian perturbiranog problema je
$f(x) + \lambda\T(h(x) - u) + \mu\T(g(x) - v)
= L(x,\lambda,\mu) - \lambda\T u - \mu\T v$. Zadnja dva člana ne ovise o $x$, pa
prolaze kroz $\inf_x$ i ostaju u dualnoj funkciji. \checkmark

\textbf{Značenje parametara.}
\begin{itemize}
  \item $v_j > 0$ znači \textbf{relaksaciju} $j$-tog ograničenja (dopuštamo
        više), $v_j < 0$ znači \textbf{pooštravanje};
  \item $u_i$ pomiče ciljanu vrijednost jednakosti.
\end{itemize}

Neka je $p\opt(u,v)$ opt.\ vrijednost primarnog problema, kao funkcija od $u$ i
$v$; $p\opt(0,0)$ je optimalna vrijednost \textbf{neperturbiranog} primarnog
problema.

\subsection{Globalna osjetljivost}
\str{32--33}

\textbf{Pretpostavimo jaku dualnost} za neperturbiran problem i neka su
$\lambda\opt, \mu\opt$ optimalne dualne varijable primarnog problema (stoga je
nužno $\mu\opt \ge 0$).

\textbf{Izvod (sa slajda 32, korak po korak):}

\begin{enumerate}
  \item Ako primijenimo \emph{slabu} dualnost na \emph{perturbiran} problem,
        imamo
        \[
        p\opt(u,v) \ge \ell(\lambda\opt,\mu\opt)
                     - (\lambda\opt)\T u - (\mu\opt)\T v
        \]
        jer je desna strana Lagrangeova dualna funkcija za perturbiran problem
        (vidjeti prethodni slajd), izračunata u jednom konkretnom dopuštenom
        paru $(\lambda\opt,\mu\opt)$ --- a svaka takva vrijednost je donja
        granica.
  \item Zbog \emph{jake} dualnosti za neperturbiran problem i iz definicije
        $\lambda\opt,\mu\opt$ imamo da je
        $\ell(\lambda\opt,\mu\opt) = p\opt(0,0)$, pa gornja nejednakost implicira
\end{enumerate}

\begin{teorem}[Globalna nejednakost osjetljivosti \str{33}]
\begin{equation}
\label{eq:globalna-osjetljivost}
p\opt(u,v) \ge p\opt(0,0) - (\lambda\opt)\T u - (\mu\opt)\T v
\end{equation}
\end{teorem}

\textbf{Implikacije (sve doslovno sa slajda 33):}

\begin{itemize}
  \item ako je $\mu_j\opt$ \textbf{velik}: $p\opt$ \textbf{znatno raste} ako
        „stisnemo” ograničenje $g_j(x)\le0$ na $g_j(x)\le v_j$ gdje je
        $v_j < 0$;
  \item ako je $\mu_j\opt$ \textbf{malen}: $p\opt$ \textbf{se ne smanjuje
        znatno} ako „relaksiramo” ograničenje $g_j(x)\le0$ na $g_j(x)\le v_j$
        gdje je $v_j > 0$;
  \item ako je $\lambda_i\opt$ \textbf{velik i pozitivan}: $p\opt$ raste znatno
        ako umjesto $h_i(x)=0$ imamo $h_i(x)=u_i$ gdje je $u_i < 0$;
  \item ako je $\lambda_i\opt$ \textbf{velik i negativan}: $p\opt$ raste znatno
        ako umjesto $h_i(x)=0$ imamo $h_i(x)=u_i$ gdje je $u_i > 0$;
  \item ako je $\lambda_i\opt$ \textbf{malen i pozitivan}: $p\opt$ se ne
        smanjuje znatno ako umjesto $h_i(x)=0$ imamo $h_i(x)=u_i$ gdje je
        $u_i > 0$;
  \item ako je $\lambda_i\opt$ \textbf{malen i negativan}: $p\opt$ se ne
        smanjuje znatno ako umjesto $h_i(x)=0$ imamo $h_i(x)=u_i$ gdje je
        $u_i < 0$.
\end{itemize}

\begin{napomena}[kako zapamtiti smjer]
Pogledajte \eqref{eq:globalna-osjetljivost}: član $-(\mu\opt)\T v$ ulazi s
\emph{minusom}. Dakle relaksacija ($v_j > 0$) \textbf{spušta} donju granicu za
$p\opt$ --- optimum se smije poboljšati (smanjiti), i to najviše za
$\mu_j\opt v_j$. Pooštravanje ($v_j < 0$) \textbf{diže} granicu --- optimum se
\emph{nužno} pogorša barem za $|\mu_j\opt v_j|$.

Uočite asimetriju: kod pooštravanja \eqref{eq:globalna-osjetljivost} daje
\textbf{jamstvo} (donja granica je čvrsta), a kod relaksacije samo
\emph{gornju ocjenu} mogućeg poboljšanja. Zato je gornja tablica formulirana s
„raste znatno” u jednom, a „ne smanjuje se znatno” u drugom smjeru.
\end{napomena}

\subsection{Lokalna osjetljivost}
\str{34}

\begin{teorem}[Lokalna osjetljivost]
Pretpostavimo jaku dualnost za neperturbiran problem i neka su
$\lambda\opt, \mu\opt$ optimalne dualne varijable primarnog problema.

Ako je $p\opt(u,v)$ diferencijabilno u $(0,0)$, tada vrijedi
\begin{equation}
\label{eq:lokalna-osjetljivost}
\lambda_i\opt = -\frac{\partial p\opt(0,0)}{\partial u_i}, \qquad
\mu_i\opt = -\frac{\partial p\opt(0,0)}{\partial v_i}
\end{equation}
\end{teorem}

\textbf{Ovo je najkorisnija formula poglavlja.} Ona kaže da Lagrangeov
množitelj \emph{jest} derivacija optimalne vrijednosti po pomaku ograničenja
--- dakle broj s izravnim inženjerskim i ekonomskim značenjem:

\begin{itemize}
  \item $\mu_j\opt$ mjeri \textbf{koliko bi se $p\opt$ popravio} kad bismo
        $j$-to ograničenje popustili za jednu jedinicu;
  \item ograničenja s $\mu_j\opt = 0$ su \textbf{neaktivna} --- njihovo
        popuštanje ne donosi ništa (usporedi uvjet komplementarnosti);
  \item u ekonomiji se $\mu_j\opt$ zove \textbf{sjenska cijena} (engl.\
        \emph{shadow price}) resursa: najviše što bismo bili spremni platiti za
        dodatnu jedinicu tog resursa.
\end{itemize}

Time se zatvara krug s fizikalnom interpretacijom iz
poglavlja~\ref{ch:uvjeti} \str{3:62--63}: množitelj je \emph{reakcija veze} ---
mjera koliko ograničenje „gura”.

\subsection{Vježbe 3, Primjer 2, podzadaci e) i f)}

\begin{zadatak}[Nastavak Zadatka~\ref{zad:primjer2}: osjetljivost]
\str{V3:40--42}

\textbf{e) Što možete reći o osjetljivosti optimalne vrijednosti na
perturbacije u ograničenjima?}

\textbf{Perturbirani problem} \str{V3:40}:
\[
\text{minimiziraj}_x \quad
\frac{1}{2}\begin{pmatrix}x_1\\x_2\end{pmatrix}\T
\underbrace{\begin{pmatrix}2&0\\0&2\end{pmatrix}}_{Q}
\begin{pmatrix}x_1\\x_2\end{pmatrix}
+ \underbrace{\begin{pmatrix}-6&-4\end{pmatrix}}_{c\T}
\begin{pmatrix}x_1\\x_2\end{pmatrix} + 13
\]
\[
\text{uz ogr.}\quad
\underbrace{\begin{pmatrix}
-1 & 0 \\ 1 & 0 \\ 0 & -1 \\ 0 & 1
\end{pmatrix}}_{G}
\begin{pmatrix}x_1\\x_2\end{pmatrix}
- \underbrace{\begin{pmatrix}0\\2\\0\\1\end{pmatrix}}_{h}
\le \begin{pmatrix}v_1\\v_2\\v_3\\v_4\end{pmatrix}
\]
Optimalno rješenje: $p\opt(v_1,v_2,v_3,v_4)$.

\textbf{Rješenje} \str{V3:41}\textbf{.} Budući da je
$\mu\opt = \begin{pmatrix}0 & 2 & 0 & 2\end{pmatrix}\T$, imamo (primjenom
\eqref{eq:lokalna-osjetljivost}):
\[
\frac{\partial p\opt(0,0,0,0)}{\partial v_1} = 0, \qquad
\frac{\partial p\opt(0,0,0,0)}{\partial v_2} = -2,
\]
\[
\frac{\partial p\opt(0,0,0,0)}{\partial v_3} = 0, \qquad
\frac{\partial p\opt(0,0,0,0)}{\partial v_4} = -2
\]

\textbf{Čitanje rezultata:}

\begin{itemize}
  \item $v_1$ i $v_3$ pripadaju \emph{neaktivnim} ograničenjima
        ($x_1 \ge 0$, $x_2 \ge 0$); njihove derivacije su \textbf{nula} ---
        pomicanje tih granica ne mijenja ništa. Logično: optimum $(2,1)$ je
        daleko od njih.
  \item $v_2$ i $v_4$ pripadaju \emph{aktivnim} ograničenjima ($x_1 \le 2$,
        $x_2 \le 1$); derivacije su $-2$. Popuštanje bilo kojeg od njih za
        malu vrijednost $\varepsilon$ smanjuje $p\opt$ za približno
        $2\varepsilon$.
\end{itemize}

\emph{Brza provjera zdravim razumom.} Ako gornju granicu na $x_1$ pomaknemo s
$2$ na $2 + \varepsilon$, novi optimum je $(2+\varepsilon,\ 1)$, pa
\[
p\opt(\varepsilon) = (2+\varepsilon-3)^2 + (1-2)^2
= (\varepsilon - 1)^2 + 1 = 2 - 2\varepsilon + \varepsilon^2
\]
Derivacija u $\varepsilon = 0$ je $-2$. \checkmark Poklapa se s
$\mu_2\opt = 2$ i formulom $\partial p\opt/\partial v_2 = -\mu_2\opt$.

\medskip
\textbf{f) Ponovite (d) i (e) s promijenjenim ograničenjima} $g_2$ i $g_4$
($g_1$ i $g_3$ ostaju ista): $g_2(x) = x_1 - 1{,}5 \le 0$,
$g_4(x) = x_2 - 1{,}5 \le 0$. \str{V3:42}

\textbf{Rezultat sa slajda:} imamo da je
$\mu\opt = \begin{pmatrix}0 & 3 & 0 & 1\end{pmatrix}\T$, pa je stoga
\[
\frac{\partial p\opt(0,0,0,0)}{\partial v_1} = 0, \qquad
\frac{\partial p\opt(0,0,0,0)}{\partial v_2} = -3,
\]
\[
\frac{\partial p\opt(0,0,0,0)}{\partial v_3} = 0, \qquad
\frac{\partial p\opt(0,0,0,0)}{\partial v_4} = -1
\]

\textbf{Provjera rezultata KKT-om} (slajd daje samo brojke; izvod niže slijedi
postupak iz podzadatka b):

\emph{Korak 1 --- novi dozvoljeni skup.} Sada je
$0 \le x_1 \le 1{,}5$, $0 \le x_2 \le 1{,}5$ --- kvadrat stranice $1{,}5$.
Ciljna točka $(3,2)$ je i dalje gore desno izvan njega, pa je najbliža točka
opet gornji desni vrh: $x\opt = (1{,}5;\ 1{,}5)$.

\emph{Korak 2 --- gradijent cilja u optimumu:}
\[
\nabla f(x\opt) = \begin{pmatrix}2\cdot1{,}5 - 6\\ 2\cdot1{,}5 - 4\end{pmatrix}
= \begin{pmatrix}-3\\-1\end{pmatrix}
\]

\emph{Korak 3 --- KKT jednadžba} (aktivni su $g_2$ i $g_4$, pa je
$\mu_1 = \mu_3 = 0$):
\[
\begin{pmatrix}-3\\-1\end{pmatrix}
+ \mu_2\begin{pmatrix}1\\0\end{pmatrix}
+ \mu_4\begin{pmatrix}0\\1\end{pmatrix} = 0
\;\Longrightarrow\;
\mu_2 = 3, \quad \mu_4 = 1
\]
Oba su $\ge 0$, dakle KKT uvjeti su zadovoljeni i
$\mu\opt = (0,3,0,1)$. \checkmark

\emph{Korak 4 --- optimalna vrijednost:}
$p\opt = (1{,}5-3)^2 + (1{,}5-2)^2 = 2{,}25 + 0{,}25 = 2{,}5$.

\textbf{Usporedba dvaju slučajeva --- što se naučilo:}

\begin{center}
\begin{tabular}{lcc}
\toprule
& \textbf{e)} $x_1\le2$, $x_2\le1$ & \textbf{f)} $x_1\le1{,}5$, $x_2\le1{,}5$ \\
\midrule
$x\opt$ & $(2;\ 1)$ & $(1{,}5;\ 1{,}5)$ \\
$p\opt$ & $2$ & $2{,}5$ \\
$\mu_2\opt$ (uz $x_1$) & $2$ & $3$ \\
$\mu_4\opt$ (uz $x_2$) & $2$ & $1$ \\
\bottomrule
\end{tabular}
\end{center}

\begin{itemize}
  \item U slučaju f) ograničenje na $x_1$ je \emph{pooštreno} (s $2$ na
        $1{,}5$), pa je njegova sjenska cijena \textbf{porasla} s $2$ na $3$ ---
        ono sada „boli” više.
  \item Ograničenje na $x_2$ je \emph{popušteno} (s $1$ na $1{,}5$), pa je
        njegova sjenska cijena \textbf{pala} s $2$ na $1$.
  \item Ukupno je $p\opt$ porastao s $2$ na $2{,}5$: gubitak zbog stiskanja
        $x_1$ nadjačao je dobitak od popuštanja $x_2$.
  \item \textbf{Praktična pouka:} ako imate ograničen budžet za popuštanje
        ograničenja, uložite ga u ono s \textbf{najvećim} $\mu\opt$. Množitelji
        vam izravno rangiraju ograničenja po tome koliko se isplati na njima
        raditi.
\end{itemize}
\end{zadatak}

% =====================================================================
\section{Sažetak poglavlja}

\textbf{Lanac ideja:}

\begin{enumerate}
  \item Za dopustiv $x$ i $\mu \ge 0$ vrijedi $L(x,\lambda,\mu) \le f(x)$ ---
        Lagrangian je \textbf{donja ograda}.
  \item Infimizacijom po $x$ dobiva se \textbf{dualna funkcija}
        $\ell(\lambda,\mu) = \inf_x L$, koja je uvijek $\le p\opt$.
  \item Maksimizacija po $(\lambda,\mu \ge 0)$ daje \textbf{dualni problem} s
        vrijednošću $d\opt$. \textbf{Slaba dualnost:} $d\opt \le p\opt$, uvijek.
  \item Dualni problem je \textbf{uvijek konveksan}, bez obzira na primarni.
  \item \textbf{Jaka dualnost} ($d\opt = p\opt$) vrijedi ako je primarni problem
        konveksan i zadovoljeni su \textbf{Slaterovi uvjeti}.
  \item Iz jake dualnosti slijede \textbf{KKT uvjeti}, i --- za konveksne
        probleme --- oni su i \textbf{dovoljni}.
  \item \textbf{Dualna dekompozicija:} dualiziranjem samo „spojnog” ograničenja
        problem se raspada na neovisne podprobleme. Primjena: formiranje
        tržišnih cijena.
  \item \textbf{Osjetljivost:} $\mu_j\opt = -\partial p\opt/\partial v_j$ ---
        množitelj je sjenska cijena ograničenja.
\end{enumerate}

\textbf{Tri stvari koje se najčešće miješaju:}

\begin{enumerate}
  \item $\ell(\lambda,\mu)$ nastaje minimizacijom po $x$ \emph{bez ograničenja}
        --- ograničenja su „ugrađena” u Lagrangian.
  \item Slaba dualnost vrijedi uvijek; \textbf{jaka} traži konveksnost
        \emph{plus} Slatera.
  \item Predznak u osjetljivosti: $\mu\opt = -\partial p\opt/\partial v$
        (s minusom), jer relaksacija ($v>0$) \emph{smanjuje} $p\opt$.
\end{enumerate}

\textbf{Terminologija (hrvatski $\leftrightarrow$ engleski):}

\begin{center}
\begin{tabular}{ll}
\toprule
\textbf{Hrvatski} & \textbf{Engleski} \\
\midrule
primarni / dualni problem & primal / dual problem \\
dualna funkcija cilja & dual function \\
dualna varijabla & dual variable \\
slaba / jaka dualnost & weak / strong duality \\
dualni jaz & duality gap \\
Slaterovi uvjeti & Slater's constraint qualification \\
dualna dekompozicija & dual decomposition \\
analiza osjetljivosti & sensitivity analysis \\
sjenska cijena & shadow price \\
marginalni (inkrementalni) trošak & marginal (incremental) cost \\
\bottomrule
\end{tabular}
\end{center}
